% How the LaTeX preview works: the article's first chapter as a LaTeX document, to open in Sagebrush's LaTeX editor.
% Written by Claude Opus 5.5, an AI model made by Anthropic, for the Sagebrush project (SageMath, Inc.).
% Same license as the rest of articles/. Numbers: bench/latex, November 4, 2026 (see 01-how-it-works.md).
\documentclass[11pt]{article}
\usepackage{amsmath,amssymb,amsthm}
\usepackage{hyperref}

\newtheorem{principle}{Principle}

\newcommand{\brush}{Sagebrush}
\newcommand{\dice}{\operatorname{Dice}}

\title{How the \LaTeX{} preview works}
\author{Claude Opus 5.5, for the \brush{} project}
\date{November 4, 2026}

\begin{document}
\maketitle

\begin{abstract}
Open a \texttt{.tex} file in \brush{} and the document appears beside its source:
headings, theorems, equations typeset by KaTeX, tables, footnotes and references.
This document is itself an example: you are reading it in that preview. Edit it,
and the preview follows.
\end{abstract}

\tableofcontents

\section{An expander, then a renderer}\label{sec:parts}

\LaTeX{} is a language of macros: \verb|\section{Intro}| means something only after
the definitions it expands to, and a document can define its own, as this one
does with \verb|\brush| and \verb|\dice|. So the preview has two parts.

\begin{description}
  \item[The expander] reads the source as \TeX{} does: category codes, control
    sequences, groups, registers, conditionals, \verb|\def| and
    \verb|\newcommand|, with bounds on work, depth and size, so a runaway macro
    stops instead of hanging the page. Where \TeX{} would typeset, it records an
    \emph{event}.
  \item[The renderer] turns the events into HTML: paragraphs, theorems, lists,
    tables, footnotes,\footnote{Like this one.} the table of contents, labels
    and references, and mathematics through KaTeX.
\end{description}

Both run in a Web Worker, so typing never waits for them. The expander is
incremental: it keeps checkpoints, resumes from the last one before an edit, and
stops once its state matches the previous run again.

\begin{principle}\label{pr:visible}
A preview that silently drops part of a document is worse than none. Every gap
should be visible, in the preview and in \texttt{sagebrush latex check}.
\end{principle}

\section{Real \TeX{} as the oracle}\label{sec:oracle}

For each document of a public corpus, the benchmark compares the words of
\TeX{}'s PDF with the words of the preview, as multisets $P$ and $Q$:
\begin{equation}\label{eq:dice}
  \dice(P, Q) = \frac{2\,|P \cap Q|}{|P| + |Q|},
  \qquad
  \text{recall} = \frac{|P \cap Q|}{|P|},
  \qquad
  \text{precision} = \frac{|P \cap Q|}{|Q|}.
\end{equation}
Mathematics and pictures differ on both sides, so $\dice = 1$ isn't expected; a
low value of \eqref{eq:dice} is the signal, and the words missing most say why.

\section{Speed}\label{sec:speed}

Medians of 5 runs on a 16-core virtual machine, with \TeX{} Live 2023 and
Chrome 155. A typed character to the preview showing it, against
\texttt{latexmk} remaking the PDF after the same edit:

\begin{center}
\begin{tabular}{lrrr}
\hline
document & \texttt{latexmk} & preview & faster \\
\hline
Beamer talk & 3.88 s & 32 ms & 120$\times$ \\
pgfplots manual, 48 pages & 21 s & 258 ms & 81$\times$ \\
amsmath user's guide & 657 ms & 74 ms & 8.9$\times$ \\
arXiv paper, 25 pages & 801 ms & 100 ms & 8.0$\times$ \\
\hline
\end{tabular}
\end{center}

\section{What would have saved time}\label{sec:lessons}

\begin{enumerate}
  \item Real \TeX{} as the oracle from the first day (Section~\ref{sec:oracle}).
  \item Measure what the reader sees: a partial paint once made a 3-second
    preview look like 240 ms.
  \item Ask \TeX{} instead of guessing.
  \item Report what you can't show, from the start (Principle~\ref{pr:visible}).
\end{enumerate}

The whole article, with the gaps and the full tables, is at
\url{https://sagebrush.space/articles/latex/}.

\end{document}
